\chapter{Welcome to Fate's Edge}
łabel{ch:intro}

\emph{A world where every choice carries weight, every spell risks backlash, and every legend is written in the shadow of consequence.}

Welcome to \textbf{Fate's Edge}, a tabletop roleplaying game where narrative drives mechanics, and every decision shapes not only your character's path−−−but the world around them. This is not a game of perfect successes or clean victories. It is a game of risk, drama, and legacy.

\section*{What Is This Game?}
∈dex{game concept}

Fate's Edge is a narrative‑first RPG where:
\begin{itemize}
ℹtem Every roll introduces potential for triumph \emph{and} complication.
ℹtem Magic is powerful−−−but dangerous.
ℹtem Choices ripple outward, shaping character arcs and the setting.
ℹtem Growth is meaningful, earned through experience spent on skills, talents, and assets.
\end{itemize}

This guide helps you build a character, understand the setting, and step into a world where your actions matter. For the complete core rules (dice resolution, Position, Difficulty Values, Story Beats, Boons, and more), refer to the \textit{Fate's Edge System Reference Document (SRD)}.

\section*{Modular by Design}
∈dex{modularity}∈dex{game structure}

One of the most important things to understand about \textit{Fate's Edge} is that it is \textbf{modular}. You are not expected to memorise every subsystem before you sit down to play. Instead, the game is designed to be discovered layer by layer, as your table’s ambition grows.

\begin{itemize}
  ℹtem \textbf{Session 1:} Use only the core resolution cycle (\autoref{ch:core−mechanics}) and maybe one regional travel generator. That is a complete, playable game.
  ℹtem \textbf{Session 5:} You will have integrated Assets, Followers, and a Patron or two. The system will feel natural.
  ℹtem \textbf{Session 20:} You will be juggling Obligation, Corruption, faction timers, and legacy planning without looking at the book.
\end{itemize}

The GM decides which subsystems to introduce and when. Your job as a player is to learn what is on your character sheet and trust that the rest will come when it is needed. This modularity means:
\begin{itemize}
  ℹtem You never need to read the entire book before playing.
  ℹtem The game grows with your table, not against it.
  ℹtem Every rule exists to serve a narrative purpose; if a subsystem is not in play, your story is not poorer for it.
\end{itemize}

So do not panic when you see a long table of Patrons or a complex ritual procedure. Start with the core loop, then add one new thing each session. The road is long; you will get there.

\section*{Core Principles}
∈dex{design philosophy}

The game is built on four key ideas:

\begin{description}
ℹtem[Narrative First] Mechanics serve the story. Rules reward descriptive play and creative problem‑solving.
ℹtem[Risk Creates Drama] Every roll carries tension. Even success may come at a cost.
ℹtem[Meaningful Growth] Experience is a currency of choice. Invest in yourself or in your influence on the world.
ℹtem[Consequences Matter] No action is free. Every choice changes the fiction.
\end{description}

\section*{Style of Play}
∈dex{tone of play}

Expect cinematic, collaborative storytelling:
\begin{itemize}
ℹtem Stories driven by character choices.
ℹtem A world that reacts to your decisions.
ℹtem Themes of legacy, sacrifice, and moral complexity.
\end{itemize}

Whether you are a lone duelist, a scheming mastermind, or a spirit‑touched outlander, your path is yours to forge.

\section*{Safety Tools: Lines, Veils & X‑Card}
łabel{sec:safety}

\begin{tcolorbox}[colback=red!5!white, colframe=red!75!black, title=\textbf{Safety Tools: Lines, Veils & X‑Card}, fonttitle=\bfseriesłarge]
Fate's Edge explores mature themes−−−conflict, loss, horror, betrayal, and difficult choices. The table should always feel safe. Use these simple tools to keep the game enjoyable for everyone.

\medskip
\noindent\textbf{X‑Card (by John Stavropoulos).} Any player (including the GM) can draw an ``X'' on a card or simply say ``X'' when a moment makes them uncomfortable. When the X‑Card is shown, the described content is \textbf{removed from the fiction}−−−no questions, no explanations, no judgments. You do not need to justify using the X‑Card. The scene is retconned or redirected immediately.

\medskip
\noindent\textbf{Lines.} A \textbf{Line} is a hard boundary: content that will not appear in the game at all. Examples: ``No sexual violence,'' ``No harm to children,'' ``No body horror.'' Before the campaign or a session, players name their Lines. The GM and group honour them completely.

\medskip
\noindent\textbf{Veils.} A \textbf{Veil} is a soft boundary: content that can happen ``off screen'' or fade to black. The event is acknowledged but not described in detail. Example: ``A romantic scene can happen, but fade to black before it becomes explicit.'' Veils allow the story to acknowledge mature themes without forcing players to roleplay through them.

\medskip
\noindent\textbf{How to Use These Tools at Your Table.}
\begin{itemize}
ℹtem Establish Lines and Veils during \textbf{Session Zero} (or before your first game). Write them down where everyone can see them.
ℹtem Keep an X‑Card visible on the table (a printed card, a drawn X, or a digital symbol). Remind players they can use it at any time, even mid‑sentence.
ℹtem When the X‑Card is used, pause, rewind if needed, and move on. Do not interrogate the player who used it.
ℹtem The GM is also a player. The GM may use the X‑Card for their own comfort.
ℹtem Respect the tools. If someone calls an X‑Card, trust that they had a good reason and support them.
\end{itemize}

\medskip
\noindent These tools are not about censorship−−−they are about \textbf{consent}. They allow the group to tell darker, more intense stories together \emph{because} everyone knows they can tap out at any time. A game that feels safe is a game that can go to the edge.

\medskip
\noindent For more detail, search online for ``TTRPG X‑Card'' or ``Lines and Veils RPG''. Many free resources are available.
\end{tcolorbox}

\section*{Session Zero: The Foundation of Your Story}
łabel{sec:session−zero−player}
∈dex{Session Zero}

Before you roll your first die, gather with your GM and fellow players for a \textbf{Session Zero}. This is where you build the shared understanding that makes Fate's Edge sing −− and where you make sure everyone is excited for the same kind of story.

\subsection*{Core Philosophy to Embrace as Players}
\begin{itemize}
    ℹtem \textbf{Optimal play = most fun.} Spending Boons, embracing failure, and courting Patron attention leads to better stories. There is no ``winning'' except an unforgettable session.
    ℹtem \textbf{You and the GM are co‑creators, not adversaries.} The GM frames pressure; you drive the fiction. No one is trying to ``beat'' anyone else.
    ℹtem \textbf{Boons are not XP tokens.} Spend them. Hoarding Boons for conversion starves the table of drama and leads to less engagement.
    ℹtem \textbf{Patrons are allies who demand payment −− embrace it.} Obligation is plot, not punishment. When a Patron calls, a story hook arrives.
\end{itemize}

\subsection*{Session Zero Checklist (Player's Edition)}
\begin{tcolorbox}[colback=green!5!white, colframe=green!75!black, title=Session Zero Checklist, fonttitle=\bfseries]
Answer these together and write down the answers where everyone can see them.
\begin{enumerate}
    ℹtem \textbf{Lines & Veils} −− What content is off‑limits (Lines) or only implied (Veils)?
    ℹtem \textbf{Tone} −− Cinematic action? Grim survival? Political intrigue? Dark horror? Mix of several?
    ℹtem \textbf{Campaign length} −− One‑shot, short arc (3−−6 sessions), or long campaign (12+ sessions)?
    ℹtem \textbf{Character hooks} −− Each player shares one goal and one fear for their character.
    ℹtem \textbf{Patron interest} −− Which Patrons (if any) are already watching the party?
    ℹtem \textbf{Starting situation} −− Where does the story begin?
    ℹtem \textbf{House rules} −− Any optional subsystems (Threat Pool, Deck of Consequences, etc.) the GM plans to use?
    ℹtem \textbf{Scheduling & attendance} −− How often do you meet? What happens if a player misses a session?
\end{enumerate}
\end{tcolorbox}

Taking thirty minutes for Session Zero will save hours of mid‑campaign friction and give you a richer, more invested table.

\section*{Guide Structure}

This Player's Guide contains:
\begin{itemize}
ℹtem \textbf{\autoref{ch:core−mechanics}} −−− Action resolution, Position, Difficulty Values, Story Beats, Boons, and basic combat.
ℹtem \textbf{\autoref{ch:advancement}} −−− Attributes, Skills, dice pools, Experience Points, and the Tier system.
ℹtem \textbf{\autoref{ch:talents}} −−− Special abilities that define your character's style.
ℹtem \textbf{\autoref{ch:magic}} −−− The five magical paths (Free Caster, Runekeeper, Invoker, Cantor, Summoner).
ℹtem \textbf{\autoref{ch:healing}} −−− How to recover Fatigue and Harm.
ℹtem \textbf{\autoref{ch:assets−followers}} −−− Off‑screen holdings and on‑screen allies.
ℹtem \textbf{\autoref{ch:currency−free−economy}} −−− No coin; use XP, Strings, and Debt Timers.
ℹtem \textbf{\autoref{ch:world−cultures}} −−− The peoples, lands, and factions of Fate's Edge.
ℹtem \textbf{\autoref{ch:world−powers−obligation}} −−− The Patrons and their Rites, Obligation, and Corruption.
ℹtem \textbf{\autoref{ch:example−concepts}} −−− Ten example characters to inspire you.
ℹtem \textbf{\autoref{ch:enhanced−play}} −−− Tactical advice, momentum banking, and optional rules.
ℹtem \textbf{\autoref{ch:downtime}} −−− What you do between adventures.
ℹtem \textbf{\autoref{ch:legacy−engine}} −−− Passing the torch to a successor.
\end{itemize}

\begin{tcolorbox}[colback=gray!5!white, colframe=gray!75!black, title=Quick Reference: Where to Find It]
∈dex{quick reference}
Use this one‑page guide to jump straight to the rules you need at the table.

\begin{tabularx}{\textwidth}{lX}
⊤rule
\textbf{What you need} & \textbf{Where to look} \\
∣rule
Core resolution (roll dice, DV, outcomes) & \autoref{ch:core−mechanics} \\
Position & Effect & \autoref{sec:position} (Core Mechanics) \\
Magic (paths, rites, backlash) & \autoref{ch:magic} \\
Travel (legs, timers, roles) & \autoref{sec:travel−framework} (World Interaction) \\
Story Beats (SB) & \autoref{sec:story−beats} (Core Mechanics) \\
Boons (earning, spending, limits) & \autoref{sec:boons} (Core Mechanics) \\
Assets & Followers & \autoref{ch:assets−followers} \\
⊥tomrule
\end{tabularx}
\end{tcolorbox}

\section*{How to Use This Book}

Read cover to cover or jump to relevant sections. Each chapter stands alone while connecting to broader themes. The SRD provides full mechanical detail; this guide focuses on character creation, setting, and player‑facing options.

\section*{Getting Started}

This is a game of bold choices and lasting consequences. Your story is written in decisions−−−not dice rolls.

Welcome to the Edge. The world is watching.

\begin{center}
\emph{What will you risk to reshape the world?}
\end{center}

\section{The Philosophy of Narrative First}
łabel{sec:narrative−first}

In Fate's Edge, the story is not a byproduct of the rules−−−it is the reason the rules exist. When you sit down to play, you are not rolling dice to see what happens; you are telling a story together, with dice as your co‑authors. This is the \textit{narrative first} philosophy.

\subsection{Why Narrative First?}

Imagine this scene:

\begin{quote}
ℹtshape Your character, a wounded duelist, faces the baron's champion in the royal arena. Your legs ache from the earlier fight, and the crowd's roar has become a physical pressure. You know one strike will end it−−−but you also know you are one misstep from death.
\end{quote}

Now consider two ways this might play out:

\begin{itemize}
ℹtem \textbf{Mechanics‑First Approach:} ``I roll Body 3 + Melee 2 to attack. I get 3 successes. The baron's champion has 2 Harm. I win.''

ℹtem \textbf{Narrative‑First Approach:} ``I lunge forward, my blade a silver arc against the red of the arena sand. But my injured leg buckles−−−\textit{I spend 1 Boon to re‑roll the die that came up 1}−−−and I twist to drive my shoulder into his chest instead. The crowd roars as we both fall to the sand, my blade at his throat.''  
\emph{GM: ``You have the advantage, but your wound has reopened. What do you do?''}
\end{itemize}

The first version tells you who won. The second tells you \textit{how} you won and what it cost you. This is the heart of narrative first: your choices matter not because of numbers on a sheet, but because of the story they create.

\subsection{Your Role in the Story}

In Fate's Edge, you are:
\begin{itemize}
ℹtem \textbf{A storyteller} −−− Your descriptions shape the world.
ℹtem \textbf{An active participant} −−− Your choices drive the narrative.
ℹtem \textbf{A co‑creator} −−− You and the GM build the story together.
\end{itemize}

This means:
\begin{itemize}
ℹtem \textbf{Flavor is free:} Describe how you fight, what your character looks like, or what the world feels like−−−no mechanical cost, only story enrichment.  
\emph{Example:} ``I parry with the grace of a dancer, my movements a memory of my mother's training, even as blood drips from my brow.''

ℹtem \textbf{Describe your action before rolling:} Do not just say ``I attack''−−−show us how. This informs the GM's response and often determines Position (Dominant, Controlled, Desperate).  
\emph{Example:} ``I use the crowd's noise as cover to slip behind him and strike his exposed back.'' (Likely \textit{Dominant Position}.)

ℹtem \textbf{Embrace failure as opportunity:} A miss is not ``you fail to hit.'' It is ``the blade glances off his armor, and you see the baron's hand reach for the hilt of his dagger.'' Failures should always move the story forward.
\end{itemize}

\subsection{The Golden Rule}

\textbf{When in doubt, make the choice that serves the story.} If you find yourself asking ``What's the optimal move?'' try reframing it as ``What would be most interesting for the story?''

\begin{tcolorbox}[colback=black!3!white, colframe=black!40!white, title=Remember, fonttitle=\bfseries]
In Fate's Edge, the dice do not determine your success or failure. \textit{You} determine the story. The dice just help decide which path you take to get there.
\end{tcolorbox}

\subsection{Flavor is Free (Even for Mundane Actions)}

\begin{tcolorbox}[colback=gray!5!white, colframe=gray!75!black, title=Flavor is Free, fonttitle=\bfseries]
∈dex{flavor}∈dex{roleplay rewards}

\textbf{Players:} Add descriptive details, cultural elements, and atmospheric touches without spending resources or requiring rolls. \textbf{This applies to all descriptive actions, not just magic.}

Want your Vhasian duelist to parry with a flourish from royal fencing schools? Go ahead. Want to describe the eerie silence of a Valewood ruin when searching for clues? Perfect. Want to narrate how you reload your crossbow with a practiced flick of the wrist? Even better.

Flavor enriches the narrative without changing mechanical outcomes. Describe your character's background, customs, or scene details. The GM should encourage this and reciprocate.

Mechanics determine success or failure, but flavor determines the story we tell about it.
\end{tcolorbox}

\section*{Narrative‑Heavy Gameplay Options (Optional)}

For groups that prefer strong narrative focus, consider these approaches:

\begin{itemize}
ℹtem \textbf{Collaborative Scene Framing:} Players may suggest scene elements (weather, NPC reactions, environmental details) that fit the fiction, with GM approval.
ℹtem \textbf{Intent‑Driven Resolution:} When success is reasonably assured, the GM may ask players to describe \emph{how} they accomplish their goal rather than rolling.
ℹtem \textbf{Flashback Declarations:} Spend 1 Boon to declare a flashback establishing past preparation, and describe the scene.
ℹtem \textbf{Descriptive Assistance:} Players can assist each other by providing vivid descriptions, granting a +1 die to the primary actor's roll.
ℹtem \textbf{Narrative Control Points:} Each player starts each session with 1 Narrative Control Point. Spend it to:
\begin{itemize}
ℹtem Introduce a minor NPC who provides useful information or assistance.
ℹtem Establish that they have a useful item on hand (within reason).
ℹtem Create a favorable environmental detail.
\end{itemize}
These points refresh each session and encourage proactive storytelling.
\end{itemize}

\section*{What's Next?}

The next chapter (\autoref{ch:core−mechanics}) covers the core mechanics that drive play. If you are eager to build a character, jump ahead to \autoref{ch:advancement} (Attributes & Skills). For an overview of the world, see \autoref{ch:world−cultures} (World Regions and Cultures).

Now close the book−−−or keep reading. The road is waiting.